% Appendix: complete specification of the RCEMIP comparison analysis.
% Written 2026-07-17 against turbulon-analysis/rcemip-deltas @ 032f0dd
% (channel pipeline) and small_pipeline.py (small-domain pipeline), with
% STEAM at turbulon-model branch flux-cascade @ c2e5516.
% \input this file after \appendix; bibliography commands (\citep) assume
% natbib-style citations already used in the manuscript.

\section{Complete specification of the LES-ensemble comparison}
\label{sect:rcemip analysis}

This appendix specifies, in full, the analysis behind the comparison of
STEAM to the RCEMIP cloud-resolving ensemble
\citep[][]{wing2018,wing2020}. The intent is that the comparison could be
reproduced exactly from this description alone. Two independent
comparisons are described: one against the elongated-channel
(\texttt{RCE\_large300}) configurations of ten models, used for the
profile, distribution, and correlation-dimension statistics, and one
against the small-domain (\texttt{RCE\_small\_les300}) configurations of
four models, used primarily for object size distributions, for which the
channel geometry is too narrow.

\subsection{Comparison data}
\label{sect:rcemip data}

\paragraph{Channel ensemble.} Ten models provide three-dimensional
snapshot output for the 300~K channel configuration: SAM, CM1, three Met
Office Unified Model variants (CASIM, RA1-T, RA1-T-nocloud), MESONH,
SCALE, UCLA-CRM, and two ICON variants (LEM and NWP). Domains are
channels of approximately $6000\times 400$~km$^2$ at 3~km horizontal
spacing ($1984$--$2048$ cells long, $128$--$144$ cells wide), with
model-specific stretched vertical grids ($74$--$98$ levels to
$33$--$39.5$~km). Each model supplies three snapshots separated by ten
days (days 80, 90, and 100 of the simulation; days 93, 103, and 113 for
the UM variants; days 80, 90, and 99 for ICON-LEM). A third ICON variant
(AES) was excluded because its archived output contains neither a height
coordinate nor pressure, so moist static energy cannot be constructed.

\paragraph{Small-domain ensemble.} Four models provide small-domain
output: SAM ($480^2$ cells), CM1 ($540^2$), DALES ($504^2$), and
ICON-LEM ($500^2$), all at 200~m horizontal spacing (domains
$96$--$108$~km on a side) with $146$--$156$ stretched vertical levels.
The final ten daily snapshots of each run are used (days 41--50), except
ICON-LEM, whose archive holds only four snapshots (days 25--28), all of
which are used.

\paragraph{Variable handling.} All humidities are converted to mixing
ratios before any further computation. Most models archive specific
quantities (mass per unit moist air, CMOR names \texttt{hus},
\texttt{clw}, \texttt{cli}); these are converted as $r = q/(1-q)$, with
the vapor-only denominator (the neglected condensate contribution to
moist-air mass changes $r$ by ${\lesssim}2\%$ where condensate is
largest). SAM archives vapor as a mixing ratio directly. Cloud liquid
and ice are taken from \texttt{clw} and \texttt{cli} (UCLA names ice
\texttt{ice}); precipitating condensate is excluded everywhere, because
several models do not archive it. SAM's small-domain (native-format)
output archives only combined non-precipitating condensate
(\texttt{QN}); it is partitioned into liquid and ice with the same
temperature ramp STEAM uses internally,
$\lambda = \mathrm{clip}\!\left[(T - 235.15)/(273.15 - 235.15),\,0,\,1\right]$,
$q_c = \lambda\,q_n$, $q_i = (1-\lambda)\,q_n$, so that the partition is
identical on both sides of that comparison.

Total water is $q_t = r_v + r_c + r_i$ and moist static energy is
$h = c_p T + g z + L_v r_v$ with $c_p = 1004$~J\,kg$^{-1}$\,K$^{-1}$,
$g = 9.81$~m\,s$^{-2}$, $L_v = 2.5\times 10^6$~J\,kg$^{-1}$ (STEAM's
constants, applied identically to every model).

\paragraph{Model-specific recoveries.} Three archives require special
handling. (i)~The ICON files carry only a level \emph{index} (top-down)
as their vertical coordinate; physical height is recovered by inverting
the frozen moist static energy field distributed alongside them,
$z = (f_{\mathrm{mse}} - 1004.64\,T - 2500800\,q_v + 333700\,q_i)/9.80665$
(the constants of that archive, used only for this recovery), averaged
horizontally at one snapshot; the level-wise horizontal standard
deviation of the recovered height is required to be below 50~m.
(ii)~MESONH's files carry no coordinates at all; the 74-level RCEMIP
standard grid archived by SCALE is assumed (the two agree level-for-level
where checkable). (iii)~UCLA-CRM stores height as the fastest dimension
and includes a below-ground ghost level, which is dropped. Where a model
archives no pressure (DALES, MESONH, SCALE, UCLA-CRM), the RCEMIP
analytic-sounding surface pressure of 1014.8~hPa is used to initialize
STEAM; otherwise the horizontal mean of the model's lowest-level pressure
is used. All reductions of single-precision fields use double-precision
accumulators.

One caveat is carried through the interpretation: MESONH's archived
channel fields are anomalously smooth relative to every other member
(near-delta anomaly distributions at 4~km), which starves the STEAM run
matched to it; whether this reflects the model or the archival processing
cannot be determined from the archive, and MESONH is therefore excluded
from the shape-sensitive statistics (anomaly distributions, correlation
dimension, size distributions).

\subsection{STEAM configuration}
\label{sect:rcemip steam config}

One deliberately untuned configuration is used throughout; nothing is
calibrated to any individual host. The spheroscale is a constant 10~m at
all heights, with turbulons isotropic below it
(\texttt{piecewise\_isotropic\_below\_spheroscale}); the horizontal Hurst
exponent is $H_h = 1/2$ and the vertical-anisotropy exponent
$H_z = 5/9$; turbulons use the Mexican-hat envelope; one scale class per
octave. The turbulon flux component is active with its defaults: flux
increments drawn from extremal ($\beta = -1$) L\'evy noise with
$\alpha = 1.8$ in four substeps per scale class, with noise amplitude
$c = 0.21$, pegged so that the realized flux intermittency is
$C_1 = 0.1$ (the calibration $C_1 = 1.674\,c^{1.8}$ inverted at
$C_1 = 0.1$).

For the channel comparison STEAM runs on a $2048\times 128$ grid at
$\Delta x = 3$~km ($6144\times 384$~km$^2$, matching the host geometry
class) with outer scale $L = 96$~km, giving five scale classes (96, 48,
24, 12, 6~km). The domain is periodic; the 384~km cross-channel extent
holds four outer-scale tiles. For the small-domain comparison STEAM runs
on a $512^2$ grid at $\Delta x = 200$~m ($102.4$~km$^2$, one outer-scale
tile) with $L = 102.4$~km, giving nine scale classes (102.4~km down to
400~m). In both cases the simulated depth is 20~km, and the vertical
grid follows the model's internal rule
$\Delta z(z) = k_z(z)/2$ with $k_z$ set by the anisotropy relation at
each scale class.

Each STEAM run is initialized from one host snapshot: the host's
horizontal-mean $h(z)$ and $q_t(z)$, interpolated to a uniform 50~m
profile grid, plus its surface pressure as above. Soft-clamp bounds are
$h \in [\min h(z) - 10\,c_p,\ \max h(z) + 10\,c_p]$ and
$q_t \in [0,\ \max(0.03,\ 1.5 \max q_t(z))]$~kg\,kg$^{-1}$. One
realization is run per host snapshot (three per model for the channel,
ten --- four for ICON-LEM --- for the small domain), with a distinct seed
per snapshot, so host and STEAM sample counts match by construction.

\subsection{Profile statistics and anomaly distributions}
\label{sect:rcemip stats}

Per-level statistics are computed on each case's native vertical grid:
horizontal mean and variance of $h$ and $q_t$, and cloud fraction
defined as the areal fraction with non-precipitating condensate
$q_c + q_i > 0.01$~g\,kg$^{-1}$. For comparison across grids all
profiles are linearly interpolated to a common 100~m grid spanning
0--20~km (levels outside a model's native range are excluded, not
extrapolated), and snapshot-averaged.

Differences are presented as $\Delta = $ STEAM $-$ host, one curve per
model. The yardstick shading is built from the \emph{pairwise}
inter-host differences: for each level, all ordered pairs
host$_i$ $-$ host$_j$ ($i \neq j$; 90 pairs for ten models) are formed
from the snapshot-averaged profiles, and the shading shows their full
range and interquartile (25--75\%) range. Pairwise differences, rather
than deviations from the ensemble mean, are the like-for-like reference
for a pairwise quantity such as $\Delta$; deviations from the mean are
smaller by roughly $\sqrt{2}$ and would overstate STEAM's exceedances.

Anomaly probability densities are computed at the model level nearest
4~km and 10~km: at each level the horizontal mean of that snapshot is
subtracted, all snapshots of a case are pooled, and the pooled sample is
histogrammed in 150 equal-width bins spanning the pooled sample's full
range (no percentile capping, which would truncate the tails that are
precisely the point of the comparison), normalized to a density.

\subsection{Optical-depth masks, correlation dimension, and size
distributions}
\label{sect:rcemip fractal}

Cloud masks are built from vertically integrated optical depth computed
from the liquid and ice water paths with the SAM bulk-optics
relationships (as implemented in the \texttt{cloudyview} package);
precipitating species are excluded as above. A column is cloudy where
$\tau > 1$. Masks are binary arrays on each case's native horizontal
grid; the three (channel) or up to ten (small-domain) masks of a case
are always analyzed together as one ensemble, never averaged
one-at-a-time.

The ensemble correlation dimension $D_2$ is computed with
\texttt{objscale} \citep{dewitt2026} using the sandbox estimator at
R\'enyi order $q = 2$ on the one-sided edge set of the masks --- the
Grassberger--Procaccia correlation integral --- with pair separations
binned between a minimum of two pixels (6~km) and a maximum of one-third
the narrow domain extent (${\approx}127$~km) for the channel, and the
package defaults (8 pixels = 1.6~km to one-third of ${\approx}96$~km =
32~km) for the small domains; sandbox centers are subsampled by a factor
of ten. The two-pixel channel minimum admits pixel-scale discretization
bias at the small end of the range and is the price of retaining one
decade of scaling in a 128-cell-wide mask; the small-domain analysis is
free of this compromise. The correlation integrals themselves are
inspected for linearity in log--log space before the fitted slope is
reported as a dimension.

Object size distributions use the truncation-corrected estimator of
\citet{dewitt2024}: objects are 4-connected components of the mask;
areas are binned logarithmically; bins in which more than half the
objects touch the domain boundary, or which hold fewer than 30 objects,
are excluded; the power-law exponent is the least-squares slope of
$\log_{10}$ counts against $\log_{10}$ area over the surviving bins. The
host channel domains are periodic, but boundary-touching objects are
treated as truncated rather than wrapped --- a conservative choice that
discards, rather than reconnects, objects crossing the seam. For the
channel geometry the surviving range is only ${\approx}1.5$ decades of
area and the ten hosts' fitted exponents disagree by a factor of four
(0.6--2.3), so the channel ensemble does not constrain the exponent at
this domain size --- consistent with the finite-domain analysis of
\citet{dewitt2024} --- and the small-domain ensemble, with objects
resolved from 400~m, is the meaningful size-distribution comparison.

\endinput
